\section{Source games and inverse matrix structure}
\label{sec:preliminaries}

This section supplies the three external ingredients of the reduction and
the precise matrix consequence that we will use.  A perfect-completeness
PCP and parallel repetition give a source game of arbitrarily small
constant value.  Grassmann expansion gives a dense constant-value slice
of a matrix table that often agrees across a rank-one step.  We prove
the extensions needed for tables with arbitrarily many possible values,
for small column spaces, and for finding a slice by random row advice.

Throughout, $\F$ denotes the field with two elements.  All vector spaces are finite-dimensional
over $\F$, and all unspecified uniform distributions are on finite sets.
For probability measures $\mu,\nu$ on a common finite set, write
$\TV(\mu,\nu)=\frac12\sum_x|\mu(x)-\nu(x)|$; in particular, every event
has probabilities differing by at most $\TV(\mu,\nu)$.

\subsection{Finite projection games and the source gap}

A \emph{finite projection game} $G$ consists of disjoint finite question
sets $U,V$, nonempty finite answer sets $L_u,R_v$, a nonempty finite set
of edge occurrences $\mathcal E$, and a probability distribution $\mu$
on $\mathcal E$.  An occurrence $e$ has endpoints $u(e),v(e)$ and a
total map $\pi_e:L_{u(e)}\to R_{v(e)}$.  Distinct occurrences may have
the same endpoints and different maps.  A deterministic strategy is a
pair of functions selecting $a(u)\in L_u$ and $b(v)\in R_v$, and
\[
 \val(G)=\max_{a,b}\Pr_{e\sim\mu}
       [\pi_e(a(u(e)))=b(v(e))].
\]
Thus each player receives only its own question, and neither receives
the edge occurrence unless that information is already determined by
its question.  Question-dependent legal answer sets are part of the
definition.  Zero-weight occurrences may be discarded.  Isolated
questions, when retained to define auxiliary tables, do not affect value.
Independent private randomness, or shared randomness independent of the
sampled edge, cannot increase value: fix complete answer functions for
all questions using the random tapes, and average their deterministic
success probabilities.

For $k\geq1$, the game $G^{\otimes k}$ samples $k$ independent
occurrences.  Each player receives its ordered tuple of questions and
returns a tuple of legal answers; acceptance requires all $k$
constraints.  The answer coordinates remain separate even when a
question is repeated.  We also set $\val(G^{\otimes0})=1$.

\begin{theorem}[Perfect-completeness PCP gap]
\label{thm:pcp}
There is a constant $\gamma_0>0$ and a deterministic polynomial-time
map taking any Boolean $3$-CNF formula $\varphi$ to a nonempty $3$-CNF
formula $\Psi$ such that every clause of $\Psi$ uses three distinct
variables.  If $\varphi$ is satisfiable, then $\Psi$ is satisfiable.
If $\varphi$ is unsatisfiable, every assignment violates at least a
$\gamma_0$ fraction of the clauses of $\Psi$.
Clause occurrences are explicitly listed, and the binary encoding
length of $\Psi$ is polynomial in that of $\varphi$.
\end{theorem}

\begin{proof}
The PCP theorem originates in the work of Arora--Safra and
Arora--Lund--Motwani--Sudan--Szegedy
\cite{AroraSafra1998,AroraLundMotwaniSudanSzegedy1998}.
We use Dinur's perfect-completeness formulation
\cite[Definition~8.1 and Theorem~8.1, February 13, 2007 author manuscript]{Dinur2007}.
It gives, in deterministic polynomial time, an explicitly listed family of
$N$ predicates of bounded arity over a fixed finite alphabet, with a
satisfying assignment in the YES case and a fixed rejected fraction
$\theta>0$ under every assignment in the NO case.
We spell out the conversion to clauses, since perfect completeness must
survive every local encoding.

Encode each alphabet symbol injectively by a fixed number of bits.  For
each predicate, require both that its queried bit blocks encode valid
symbols and that the decoded symbols satisfy the predicate.  A satisfying
symbol assignment satisfies all these Boolean predicates.  Conversely,
decode every invalid block to a fixed default symbol.  Whenever an original
predicate rejects this decoded assignment, the corresponding Boolean
predicate rejects the bit assignment: either a queried block is invalid,
or all are valid and the original rejection is unchanged.  Thus every
Boolean assignment still violates at least $\theta N$ predicates on a
NO input.  After identifying repeated queried variables, each predicate
involves at most a fixed number $k_0$ of distinct bits.

For each rejecting assignment of a Boolean predicate, form the clause that
excludes precisely that assignment.  The conjunction of these at most
$2^{k_0}$ clauses is exactly the predicate.  Each clause has width at most
$k_0$; a rejecting zero-bit predicate contributes the empty clause.
For a clause $\ell_1\vee\cdots\vee\ell_k$ with $k\ge4$, introduce fresh
variables $y_1,\ldots,y_{k-3}$ and replace it by
\[
 (\ell_1\vee\ell_2\vee y_1),\qquad
 (\neg y_j\vee\ell_{j+2}\vee y_{j+1})\ (1\le j\le k-4),\qquad
 (\neg y_{k-3}\vee\ell_{k-1}\vee\ell_k).
\]
If all $\ell_h$ are false, the first clause forces $y_1=1$, each middle
clause propagates this value, and the last clause fails.  If the original
clause is true, setting $y_j=1$ exactly when
$\ell_1,\ldots,\ell_{j+1}$ are all false satisfies every displayed clause.
Use disjoint auxiliary variables for every such replacement.  Clauses of
width at most three are retained.  Consequently an accepting predicate
admits a satisfying extension, whereas a rejecting predicate forces a
violated clause for every extension.  There is a fixed bound $B\ge1$ on
the number of clauses per predicate.  The resulting formula has at most
$BN$ clauses and at least $\theta N$ violated clauses under every assignment
on a NO input, so its gap is at least $\widehat\gamma=\theta/B>0$.

Here is the further normalization, including its gap loss.  Remove
tautological clauses and simplify repeated literals in each remaining
clause.  A resulting clause $C$ on $k\leq3$ distinct variables is
replaced by all $2^{3-k}$ clauses obtained by adjoining every sign
pattern on $3-k$ fresh variables, used only for this occurrence.  If
$C$ is true, every extension is true.  If $C$ is false, exactly one
extension is false for every assignment to the fresh variables.
This also handles an empty clause, using all eight sign patterns on
three fresh variables.  Consequently every violated original clause
still contributes a violated clause, while the number of clauses
increases by at most eight.  The gap is therefore at least
$\widehat\gamma/8$, and completeness remains one.  If normalization
leaves no clauses, use a single satisfiable clause on three fresh
variables; this case cannot occur for a NO output of the gap reduction.
An input with no clauses can likewise be sent directly to this fixed
satisfiable formula.  All steps are deterministic and polynomial.
\end{proof}

\begin{theorem}[Weighted projection-game repetition]
\label{thm:repetition}
Let $G$ be any finite projection game with the conventions above.
For $0<g<1$ and every integer $k\geq1$,
\begin{equation}
 \val(G)\leq1-g
 \quad\Longrightarrow\quad
 \val(G^{\otimes k})\leq(1-g^2/16)^k.
 \label{eq:prelim-repetition}
\end{equation}
The bound is uniform in the question sets, answer sets, and probability
distribution on edge occurrences.
\end{theorem}

Raz proved exponential decay under parallel repetition for general two-player
games~\cite{Raz1998}; Holenstein simplified that proof~\cite{Holenstein2009},
and Rao obtained an alphabet-independent bound for projection games
\cite{Rao2011}.  The precise estimate above is the parallel repetition theorem
of Dinur and Steurer
\cite[Corollary~1.2]{DinurSteurer2014}.  Their Sections~2.1--2.2
allow nonnegative weights and parallel edges and use independent
products of questions.  Interchanging the players matches our
projection direction.  To express question-dependent answer sets in
their common-alphabet convention, embed all legal answers in one
finite alphabet and give illegal answers no accepting pairs.  Such
partial projection constraints are allowed there, and legal defaults
show that this extension does not change value.  Thus the stated
version applies directly to our games.

We now specify the source questions, because their exact local domains
will matter in the construction.  For a normalized clause occurrence
$C$, let $\mathcal A_C\subseteq\F^3$ be its seven satisfying
assignments, recorded as \emph{variable bits}, including the signs of
the clause only in the definition of $\mathcal A_C$.  The game
$G_\Psi$ samples a uniform clause occurrence $C$ and a uniform
position $h\in\{1,2,3\}$.  The left question is $C$, with answer
domain $\mathcal A_C$; the right question is the variable ID at
position $h$, with answer domain $\F$; and the projection selects
the $h$th bit.  The occurrence and selected position remain with the
verifier.

A satisfying assignment to $\Psi$ gives value one.  Conversely,
fix any right labeling, hence an assignment to the variables of
$\Psi$.  At every clause it violates, any legal left answer differs
from that assignment in at least one position.  Therefore
\begin{equation}
 \varphi\text{ unsatisfiable}
 \quad\Longrightarrow\quad
 \val(G_\Psi)\leq1-\gamma_0/3.
 \label{eq:prelim-clause-gap}
\end{equation}
For any required constant $\sigma>0$, choose a fixed integer $t$ with
$(1-\gamma_0^2/144)^t\leq\sigma$ and take
\begin{equation}
 G_{\mathrm{src}}=G_\Psi^{\otimes t}.
 \label{eq:prelim-source-game}
\end{equation}
Theorem~\ref{thm:repetition} gives source soundness at most $\sigma$,
with perfect completeness.  Its left question is an ordered tuple of
clause occurrences, its answer domain is the product of their
$\mathcal A_C$, its right question is the corresponding tuple of
variable IDs, and its right answer domain is always $\F^t$.
In particular, a repeated variable ID in a right question imposes no
equality condition on the corresponding answer coordinates.
Both answer-domain sizes, $7^t$ and $2^t$, are constants, and the
explicit game has polynomial size because $t$ is fixed.

Each $\mathcal A_C$ projects onto both bits at any single position and
onto all four bit pairs at any two distinct positions.  Indeed, a
clause excludes just one of the eight triples, and any prescribed pair
has two extensions.  Taking products proves that every source
projection is surjective.  The pair-surjectivity observation will also
make the support reduction of function queries unambiguous.

\subsection{Matrix slices and the Grassmann input}

Let $H$ be a finite vector space, put $K=\F^\ell$, and identify
$K\otimes H$ with $\Hom(H^*,K)$.  Thus a matrix $M$ is a list of
$\ell$ rows in $H$, and $Mq\in K$ evaluates those rows at
$q\in H^*$.  For a linear map $A:K\to\F^a$, the notation $AM$
denotes its $a$ specified combinations of rows.  A \emph{row and
column slice} is a nonempty affine set
\begin{equation}
 \mathcal S=
 \{M\in K\otimes H: AM=U,\ Mq_b=t_b\ (1\leq b\leq c)\},
 \label{eq:prelim-slice}
\end{equation}
where $U\in\F^a\otimes H$, $q_b\in H^*$, and $t_b\in K$.
There is no independence requirement on the rows of $A$ or on the
$q_b$.  We count specified row combinations and column values, rather
than their ranks; redundant equations can always be removed.  A slice
with no column constraints is a \emph{row fiber}.

For $a\in K$ and $h\in H$, write $ah=a\otimes h$.
The matrix step used below is
\[
 M\longmapsto M+ah,
 \qquad M\text{ uniform},\quad
 a\text{ uniform in }K\setminus\{0\},\quad h\text{ uniform in }H,
\]
with the three choices independent.  This step is stationary: both
endpoints are uniform.  The possibility $h=0$ is retained.

We recall the precise structural theorem behind the matrix argument.
The Grassmann graph on an $N$-dimensional binary vector space has
the $\ell$-dimensional subspaces as its vertices; distinct vertices
are adjacent when their intersection has dimension $\ell-1$.
For a nonempty vertex set $S$, its \emph{retention} is the probability
that a uniform vertex of $S$ has a uniformly chosen neighbor in $S$.

\begin{theorem}[Grassmann expansion consequence]
\label{thm:prelim-grassmann}
For each $0<\zeta<1$, there are $\alpha>0$, an integer $r\geq1$,
and $\ell_0$ such that, for every $\ell\geq\ell_0$ and then every
sufficiently large $N$, the following holds.  Every nonempty set of vertices
with retention at least $\zeta$ has density at least $\alpha$ in an
interval
\[
 \{L:A_0\subseteq L\subseteq B_0,\ \dim L=\ell\},
 \qquad \dim A_0+\operatorname{codim}B_0\leq r.
\]
The constants $\alpha,r,\ell_0$ depend only on $\zeta$; the
threshold for $N$ may also depend on $\ell$.
\end{theorem}

For sets of density at most $1/2$, this is Theorem~1.12 of
Khot--Minzer--Safra in ECCC TR18-006, revision~2
\cite{KhotMinzerSafra2018}, with expansion written as one minus retention;
Definitions~1.8 and~1.10 give the expansion and interval conventions.
If the set has density greater than $1/2$, take $A_0=\{0\}$ and
$B_0=\F^N$.  This whole interval has co-order zero and density greater
than $1/2$, so replacing the cited density constant by its minimum with
$1/2$ proves the statement in that case too.
The order of the dimension quantifiers is important.  We
next give the matrix-chart argument, as in
\cite[Section~3]{BarakKothariSteurer2018}, and then remove the
column-dimension requirement explicitly.

\begin{lemma}[Matrix set extraction in large dimension]
\label{lem:prelim-matrix-set}
For every $0<\eta<1$ there are $\alpha>0$, $r\geq1$, and
$\ell_0$, depending only on $\eta$, such that for each
$\ell\geq\ell_0$ and all sufficiently large $m$ the following holds.
If a nonempty set $S\subseteq\Mat_{\ell,m}(\F)$ satisfies
\[
 \Pr_{M\text{ uniform in }S,\ a\ne0,\ h}
       [M+ah^{\mathsf t}\in S]\geq\eta,
\]
where $a\in\F^\ell\setminus\{0\}$ and $h\in\F^m$ are
independent uniform vectors, then $S$ has density at least $\alpha$
in a nonempty row and column slice with at most $r$ equations of each
kind.
\end{lemma}

\begin{proof}
Increase the lower threshold on $m$ so that $2^{-m}\leq\eta/2$.
The step has $h=0$ with probability $2^{-m}$.  Its retention
conditional on $h\ne0$ is therefore at least
$(\eta-2^{-m})/(1-2^{-m})\geq\eta/2$.

Associate to $M$ its graph subspace
\[
 L_M=\{(x,M^{\mathsf t}x):x\in\F^\ell\}
       \subseteq\F^\ell\oplus\F^m.
\]
These subspaces form the chart $\mathcal C$ on which projection to the
first summand is invertible, and $M\mapsto L_M$ is a bijection onto
that chart.  The identity
$\dim(L_M\cap L_{M'})=\ell-\rank(M-M')$ shows that chart neighbors
are precisely the nonzero rank-one changes.  Over $\F$, each nonzero
rank-one matrix has exactly one factorization $ah^{\mathsf t}$ with
both factors nonzero.  The degree within the chart is thus
$(2^\ell-1)(2^m-1)$.

For comparison, a vertex of the full Grassmann graph has
$(2^\ell-1)(2^{m+1}-2)$ neighbors: choose a hyperplane in the vertex,
then a one-dimensional extension in its quotient, excluding the
original vertex.  Exactly half the neighbors of every chart vertex
therefore remain in the chart.  The set
$\widetilde S=\{L_M:M\in S\}$ has full Grassmann retention at least
$\eta/4$.  Apply Theorem~\ref{thm:prelim-grassmann} with
$\zeta=\eta/4$.  It supplies an interval $[A_0,B_0]$ of bounded
co-order on which $\widetilde S$ has density at least $\alpha$.
Its intersection with $\mathcal C$ is nonempty, and restricting the
denominator to this intersection cannot decrease that density,
since $\widetilde S\subseteq\mathcal C$.

It remains to identify the intersection as a slice.  A basis of
$A_0$ consists of vectors $(d_b,s_b)$, and its containment in $L_M$
is exactly $d_b^{\mathsf t}M=s_b^{\mathsf t}$ for every basis vector.
Using the standard dot product, a basis $(t_b,q_b)$ of $B_0^\perp$
gives the equations $Mq_b=t_b$ for containment $L_M\subseteq B_0$.
Thus there are at most $r$ row equations and at most $r$ column
equations.  The Grassmann theorem fixes $\alpha,r,\ell_0$ using
only $\eta$; after fixing $\ell$, its ambient-dimension threshold
is met by taking $m$ sufficiently large.
\end{proof}

Our eventual matrices may have very small column spaces, including
quotient spaces of dimension zero.  Adding unused columns preserves
the matrix test, but a dense slice in the enlarged space can mix real
and dummy columns.  The next lemma describes its projection exactly.

\begin{lemma}[Eliminating dummy columns]
\label{lem:padding}
Let $H,D,K$ be finite binary vector spaces, and write a matrix on
$H\oplus D$ as $(X,Y)$, with $X\in\Hom(H^*,K)$ and
$Y\in\Hom(D^*,K)$.  Suppose the nonempty affine slice
\begin{equation}
 \widetilde{\mathcal S}=
 \{(X,Y):AX=U_H,\ AY=U_D,\
       Xq_b+Yd_b=t_b\ (1\leq b\leq c)\}
 \label{eq:prelim-padded-slice}
\end{equation}
has $A:K\to\F^a$, $q_b\in H^*$, and $d_b\in D^*$.
Its image $\mathcal S$ under $(X,Y)\mapsto X$ is a nonempty row
and column slice with the same $a$ row combinations and at most $c$
column values.  Every $X\in\mathcal S$ has equally many lifts in
$\widetilde{\mathcal S}$.  In particular, a uniform point of the
padded slice projects to a uniform point of $\mathcal S$.
\end{lemma}

\begin{proof}
Put $W=\ker A$.  Encode all dependencies among the column equations
by the linear maps
\[
 q:\F^c\to H^*,\quad q(\lambda)=\sum_b\lambda_bq_b,
 \qquad
 d:\F^c\to D^*,\quad d(\lambda)=\sum_b\lambda_bd_b,
 \qquad
 t(\lambda)=\sum_b\lambda_bt_b.
\]
We claim that the image consists exactly of the matrices $X$ obeying
\begin{equation}
 AX=U_H,\qquad Xq(\lambda)=t(\lambda)
                         \quad(\lambda\in\ker d).
 \label{eq:prelim-projected-slice}
\end{equation}
A basis of $\ker d$ gives at most $c$ column equations, whether or
not their real arguments are independent.  Necessity follows by
taking linear combinations of the original column equations whose
dummy arguments vanish.

For sufficiency, fix one point $(X_*,Y_*)$ of the original nonempty
slice, and let $X$ satisfy Equation~\eqref{eq:prelim-projected-slice}.
The row equations give $A(X_*-X)=0$.  Define on $\im d$ the candidate
linear map
\begin{equation}
 L_0(d(\lambda))=(X_*-X)q(\lambda).
 \label{eq:prelim-dummy-extension}
\end{equation}
It is well-defined: if $d(\lambda)=0$, the projected equations for
$X$ and $X_*$ give $(X_*-X)q(\lambda)=0$.  Its values lie in $W$
because $A(X_*-X)=0$.  Extend $L_0$ arbitrarily to a linear map
$L:D^*\to W$, and put $Y=Y_*+L$.  Then $AY=U_D$, and
\[
 Xq(\lambda)+Yd(\lambda)
 =Xq(\lambda)+Y_*d(\lambda)+(X_*-X)q(\lambda)
 =t(\lambda)
\]
for every $\lambda$.  This gives a lift.  It also shows explicitly
why the row compatibility and all dependent column equations suffice.

Every lift is obtained this way, since its difference from $Y_*$
takes values in $W$ and must restrict to $L_0$ on $\im d$.  The
number of such extensions is
\[
 2^{\dim W\, (\dim D-\rank d)},
\]
independent of $X$.  Thus all projection fibers have the same size,
which proves the uniform-measure assertion.  The argument permits
zero dimensions and rank-deficient $A$ throughout.
\end{proof}

\subsection{Partial tables and many row fibers}

A partial table with finite range $\mathcal Y$ is a function
$f:K\otimes H\to\mathcal Y\cup\{\bot\}$, where $\bot\notin\mathcal Y$
means undefined.  A \emph{defined equality} requires both endpoints to
have the same value in $\mathcal Y$; two undefined values never count.

\begin{theorem}[Inverse theorem for a finite-range partial table]
\label{thm:inverse}
For every $0<\eta<1$, there are $\alpha>0$, integers $r\geq1$ and
$\ell_0$, depending only on $\eta$, with the following property.
For every $\ell\geq\ell_0$, every finite binary vector space $H$,
every finite set $\mathcal Y$, and every partial table
$f:K\otimes H\to\mathcal Y\cup\{\bot\}$, where $K=\F^\ell$,
if
\begin{equation}
 \Pr_{M,a,h}[f(M)=f(M+ah)\ne\bot]\geq\eta,
 \label{eq:prelim-defined-equality}
\end{equation}
for independent uniform $M$, $a\in K\setminus\{0\}$, and $h\in H$,
then some $y\in\mathcal Y$ has density at least $\alpha$ on a
nonempty slice with at most $r$ row combinations and at most $r$
column values.  The constants are independent of both $\dim H$ and
$|\mathcal Y|$.
\end{theorem}

\begin{proof}
Use the constants of Lemma~\ref{lem:prelim-matrix-set} for $\eta$.
Add a dummy space $D$ so that $\dim(H\oplus D)$ meets its column
threshold, and define $\widetilde f(X,Y)=f(X)$.  A uniform padded
matrix projects to a uniform real matrix; likewise, a uniform step
factor in $H\oplus D$ projects to a uniform factor in $H$.
Thus the defined-equality probability is unchanged, even when $H=0$.

For a defined value $y$, let $S_y=\widetilde f^{-1}(y)$, let
$\mu_y$ be its uniform matrix measure, and let $\theta_y$ be its
retention under the padded matrix step.  Empty fibers can be omitted.
Partitioning according to the first value gives exactly
\begin{equation}
 \Pr[\widetilde f(\widetilde M)
       =\widetilde f(\widetilde M+a\widetilde h)\ne\bot]
 =\sum_{y\in\mathcal Y}\mu_y\theta_y.
 \label{eq:prelim-fiber-average}
\end{equation}
Since $\sum_y\mu_y\leq1$, a value $y$ has $\theta_y\geq\eta$.
This averaging, also used in
\cite[Lemma~2.3]{BarakKothariSteurer2018}, incurs no factor for
the number of fibers.  Lemma~\ref{lem:prelim-matrix-set} supplies a
padded slice on which that fixed value has density at least $\alpha$.
Project the slice using Lemma~\ref{lem:padding}.  Its uniform measure
projects uniformly, and $\widetilde f$ depends only on the projected
matrix.  The same value therefore has the same density on the
resulting slice in $K\otimes H$, with no increase in either count of
constraints.
\end{proof}

We will need random row advice to encounter such a slice, rather than
merely the existence of one slice somewhere in the table.  The next
erasure argument supplies exactly this averaged guarantee.

\begin{lemma}[Many good row fibers]
\label{lem:many-fibers}
Fix $0<\eta<1$, and let $\alpha,r,\ell_0$ be constants from
Theorem~\ref{thm:inverse} with threshold $\eta/2$.
Suppose $\ell\geq\ell_0$ and a partial table on $K\otimes H$ has
defined-equality probability at least $\eta$.  For each linear map
$A:K\to\F^r$, call a nonempty row fiber $\{M:AM=U\}$
\emph{good} if there are at most $r$ column specifications and a
defined value $y$ such that the resulting slice is nonempty and has
$f=y$ on at least an $\alpha$ fraction of its uniform points.
Then
\begin{equation}
 \Pr_M[\text{$M$ lies in a good row fiber for some $A$}]
 \geq\eta/4.
 \label{eq:prelim-good-union}
\end{equation}
Consequently, for independent uniform $A\in\Hom(K,\F^r)$ and
$M\in K\otimes H$,
\begin{equation}
 \Pr_{A,M}[\text{the fiber of $AM$ for $A$ is good}]
 \geq (\eta/4)2^{-r\ell}.
 \label{eq:prelim-many-fibers}
\end{equation}
The witnesses can be chosen deterministically from $f,A,U$ alone.
For every good row fiber, any chosen nonempty slice defined by at most
$r$ column values has conditional uniform probability at least
$2^{-r\ell}$ inside that row fiber.
\end{lemma}

\begin{proof}
Let $T$ be the union of all good row fibers, over all the
$2^{r\ell}$ possible row maps.  Suppose its uniform measure were
less than $\eta/4$.  Erase the table on $T$, obtaining $f_0$.
Stationarity of the matrix step gives
\[
 \Pr[f_0(M)=f_0(M+ah)\ne\bot]
 \geq \Pr[f(M)=f(M+ah)\ne\bot]
          -\Pr[M\in T]-\Pr[M+ah\in T]
 >\eta/2.
\]
Theorem~\ref{thm:inverse} gives a defined value of $f_0$ of density
at least $\alpha$ on a nonempty slice with at most $r$ row and $r$
column specifications.  Pad its row map by zero output coordinates
to make its codomain $\F^r$.  Because $f_0$ is a restriction of
$f$, this slice certifies that its row fiber is good for $f$.
That entire fiber belongs to $T$, where $f_0$ is undefined, a
contradiction.  This proves Equation~\eqref{eq:prelim-good-union}.

For each $M\in T$, at least one of the $2^{r\ell}$ maps has a
good fiber through $M$.  Averaging the number of such maps gives
Equation~\eqref{eq:prelim-many-fibers}.  No statement for every
particular row map is needed or asserted.

All vector spaces and value sets involved are finite.  After fixing
orders on them, select the first witnessing list of columns, values,
and $y$ for each $(f,A,U)$.  This choice uses the row advice and
the table, without inspecting any further part of $M$.
Finally, the column equations restricted to a nonempty row fiber form
an affine linear system with at most $r\ell$ scalar equations.
Whenever it is consistent its solution set has relative size
$2^{-b}$, where $b\leq r\ell$ is the rank of that system.  This
proves the stated conditional probability, including all dependent
and zero-rank cases.
\end{proof}

All three matrix conclusions apply to quotient spaces without any
change of constants.  In particular, when $H$ is replaced by $H/B$,
the column functionals produced above belong to $(H/B)^*$; their
pullbacks to $H$ annihilate $B$.  The proof adds dummy columns to the
quotient itself and removes them by Lemma~\ref{lem:padding}, so no
functional outside this annihilator is introduced.
